\chapter{Wymagania i~narzędzia}
\label{ch:wymagania-i-narzedzia}

\section{Założenia projektowe}
\label{sec:zalozenia_projektowe}

\subsection{Sposób przetwarzania danych}
Podstawowym założeniem jest przetwarzanie danych w~całym torze na stałej szynie danych o~szerokości wynoszącej 64 bity. Dane są przekazywane pomiędzy blokami jako kolejne słowa 64-bitowe, taktowane wspólnym sygnałem zegarowym \texttt{clk}. W konsekwencji:
\begin{itemize}
    \item granice istotnych fragmentów danych nie muszą pokrywać się z~granicami słowa 64-bit,
    \item w~przypadku długości niebędącej wielokrotnością 64 bitów konieczne jest stosowanie dopełniania oraz przekazywanie informacji o~ważności danych.
\end{itemize}

\subsection{Model pracy}
Aktualna wersja systemu obsługuje dokładnie jeden blok transportowy w~ramach jednego cyklu pracy. Po zakończeniu przetwarzania TB nie przewidziano automatycznego powrotu do stanu gotowości. Aby rozpocząć przetwarzanie następnego bloku danych wymagane jest wykonanie resetu układu.

Założenie to upraszcza sterowanie modułami, zamykanie przetwarzania i~obsługę buforów danych, kosztem braku pracy w~trybie ciągłym. Rozszerzenie do obsługi transmisji ciągłej stanowi potencjalny kierunek rozwoju.

\subsection{Sterowanie przepływem danych}
W projekcie do kontroli przepływu danych zaimplementowano mechanizm opierający się na sygnale "valid" oraz "last". Mechanizm ten ideą przypomina protokół "valid-ready". 
\begin{itemize}
    \item źródło danych sygnalizuje ważność słowa sygnałem \texttt{valid},
    \item odbiornik nie posiada sygnału informującego o~gotowości do przyjęcia danych (brak \texttt{ready}),
    \item odbiornik ma obowiązek przyjąć dane w~każdym cyklu, w~którym \texttt{valid=1}.
\end{itemize}

Przyjęcie takiego interfejsu upraszcza połączenia pomiędzy blokami oraz ogranicza liczbę przypadków brzegowych podczas pierwszej implementacji. Jednocześnie nakłada wymagania na strukturę toru. Moduły muszą przekazywać dane z~wejścia na wyjście w~cyklu ich wystąpienia albo posiadać odpowiednie bufory umożliwiające przyjęcie danych w~każdym cyklu.


\subsection{Sposób podawania danych wejściowych}
Nie narzuca się stałego tempa dostarczania danych wejściowych. System nadrzędny może:
\begin{itemize}
    \item podawać dane w~sposób ciągły -- do 1 słowa 64-bitowego na 1 cykl sygnału zegarowego,
    \item wprowadzać dowolne przerwy pomiędzy kolejnymi słowami.
\end{itemize}

Słowo jest uznawane za poprawne i~podlegające przetwarzaniu wyłącznie w~cyklu, w~którym \texttt{in\_valid=1}. Zakończenie TB jest sygnalizowane przez \texttt{in\_last=1} jednocześnie z~\texttt{in\_valid=1} na ostatnim słowie TB.

Wprowadzenie pełnego mechanizmu handshake (np. zgodnego z~AXI4-Stream) jest traktowane jako możliwy etap rozbudowy, lecz nie jest elementem aktualnej wersji projektu.

\subsection{Elastyczność kodowania LDPC i~obsługa różnych TB}
Projekt zakłada obsługę dowolnego bloku transportowego, tj. poprawne przetworzenie i~zakodowanie danych niezależnie od długości TB oraz konfiguracji wynikającej z~reguł doboru parametrów kodowania LDPC. Oznacza to, że tor przetwarzania powinien działać poprawnie dla różnych przypadków, w~których mogą zmieniać się m.in.:
\begin{itemize}
    \item wybór grafu bazowego (BG),
    \item rozmiar podmacierzy,
    \item docelowe długości bloków używane w~kodowaniu,
    \item liczba bitów dopełnienia oraz wynikające z~niej wyrównanie danych,
    \item liczba CB powstałych w~wyniku segmentacji TB.
\end{itemize}

W konsekwencji architektura musi umożliwiać przetwarzanie bloków o~zmiennej długości oraz generowanie poprawnego strumienia wyjściowego przy zachowaniu stałej szerokości słowa 64-bit.

\subsection{Wymagania}

\begin{itemize}
    \item \textbf{Deterministyczność działania:} dla identycznych danych wejściowych oraz identycznych metadanych sterujących system powinien generować identyczny strumień wyjściowy.
    \item \textbf{Weryfikacja przez symulację:} projekt powinien umożliwiać weryfikację w~symulacji poprzez porównanie wyników z~wektorami referencyjnymi w~sposób bitowo-dokładny.
    \item \textbf{Syntezowalność:} projekt powinien dać się poprawnie zsyntetyzować dla założonego układu FPGA, tak aby raporty syntezy pozwalały ocenić wykorzystanie zasobów oraz spełnienie ograniczeń czasowych.

\end{itemize}


\section{Interfejsy i~format danych}
\label{sec:interfejsy_i_format_danych}

W projekcie zastosowano prosty interfejs strumieniowy oparty o~sygnały \texttt{valid/last} oraz równoległą magistralę danych 64-bit. Interfejs ten jest używany zarówno na wejściu całego toru, jak i~na wyjściu. Wewnątrz projektu poszczególne bloki mogą dodatkowo wymieniać metadane sterujące przetwarzaniem.

\subsection{Interfejs wejściowy}
\label{subsec:interfejs_wej}

Wejście projektu przyjmuje TB w~postaci kolejnych słów 64-bitowych:
\begin{itemize}
    \item \texttt{in\_chunk[63:0]} -- dane wejściowe,
    \item \texttt{in\_valid} -- informacja o~tym, że \texttt{in\_chunk} jest ważne w~bieżącym cyklu,
    \item \texttt{in\_last} -- znacznik ostatniego słowa TB,
    \item \texttt{tb\_len} -- długość TB wyrażona w~ilości bitów.
\end{itemize}

Bity w~strumieniu są interpretowane w~kolejności \textbf{MSB-first}, tj. \texttt{in\_chunk[63]} odpowiada pierwszemu bitowi danego słowa w~sensie strumieniowym, a~\texttt{in\_chunk[0]} -- ostatniemu.

Interfejs wejściowy przedstawiono na rys. \ref{fig:pdschtx_if_in}.

\begin{figure}[!ht]
\centering
\begin{minipage}{0.95\linewidth}
\begin{lstlisting}
    ...
    input  wire         clk,
    input  wire         rst,
    input  wire [63:0]  in_chunk,
    input  wire         in_valid,
    input  wire         in_last,
    input  wire [15:0]  tb_len,
    ...
);
\end{lstlisting}
\end{minipage}
\caption{Interfejs wejściowy układu.}
\label{fig:pdschtx_if_in}
\end{figure}





\subsection{Interfejs wyjściowy}
\label{subsec:interfejs_wyj}

Wyjście projektu generuje strumień danych w~postaci słów 64-bitowych:
\begin{itemize}
    \item \texttt{out\_chunk[63:0]} -- dane wyjściowe,
    \item \texttt{out\_valid} -- informacja, że \texttt{out\_chunk} jest ważne w~bieżącym cyklu,
    \item \texttt{out\_last} -- znacznik ostatniego słowa,
    \item \texttt{out\_pad\_left[5:0]} -- liczba bitów do pominięcia od strony MSB w~słowie wyjściowym,
    \item \texttt{out\_pad\_right[5:0]} -- liczba bitów do pominięcia od strony LSB w~słowie wyjściowym.
\end{itemize}

Słowo wyjściowe jest uznawane za ważne wyłącznie w~cyklu, w~którym \texttt{out\_valid=1}. Zakończenie przetwarzania całego TB jest jednoznacznie określone przez \texttt{out\_last=1} podane wraz z~ostatnim słowem wyjściowym.

Analogicznie do wejścia, bity w~strumieniu wyjściowym są interpretowane w~kolejności \textbf{MSB-first}, tj. \texttt{out\_chunk[63]} odpowiada pierwszemu bitowi danego słowa w~sensie strumieniowym, a~\texttt{out\_chunk[0]} -- ostatniemu.

Sygnały \texttt{out\_pad\_left} oraz \texttt{out\_pad\_right} stanowią metadane umożliwiające jednoznaczną interpretację strumienia w~przypadku, gdy liczba ważnych bitów w~słowie jest mniejsza niż 64. Informują one, ile bitów w~słowach należy pominąć odpowiednio od strony MSB oraz LSB.

Interfejs wyjściowy przedstawiono na rys. \ref{fig:pdschtx_if_out}.

\begin{figure}[!ht]
\centering
\begin{minipage}{0.95\linewidth}
\begin{lstlisting}
    ...
    output wire [63:0]  out_chunk,
    output wire         out_valid,
    output wire         out_last,
    output wire [5:0]   out_pad_left,
    output wire [5:0]   out_pad_right
);
\end{lstlisting}
\end{minipage}
\caption{Interfejs wyjściowy układu.}
\label{fig:pdschtx_if_out}
\end{figure}


\section{Środowisko projektowe i~narzędzia}
\label{sec:srodowisko_i_narzedzia}

Do realizacji projektu wykorzystano zestaw narzędzi obejmujący środowisko do syntezy i~implementacji układów FPGA, symulator HDL oraz środowisko do przygotowania i~analizy danych referencyjnych.

\subsection{Quartus Prime Lite Edition}
Podstawowym narzędziem do syntezy, implementacji oraz analizy wykorzystania zasobów i~parametrów czasowych jest Quartus Prime. Narzędzie to zostało użyte do:
\begin{itemize}
    \item syntezy logiki opisanej w~języku Verilog,
    \item implementacji (fitter) dla wybranego układu FPGA,
    \item generowania raportów wykorzystania zasobów,
    \item analizy czasowej (timing) oraz wyznaczenia maksymalnej częstotliwości pracy.
\end{itemize}

\subsection{Questa Intel Starter FPGA Edition}
Symulacje funkcjonalne projektu oraz uruchamianie środowiska testowego zrealizowano z~wykorzystaniem symulatora Questa. Narzędzie to zostało użyte do:
\begin{itemize}
    \item kompilacji i~uruchamiania \english{testbench'y} HDL napisanych w~języku System Verilog,
    \item analizy przebiegów czasowych oraz zawartości rejestrów w~celu debugowania,
    \item weryfikacji poprawności działania poprzez porównanie wyników z~danymi referencyjnymi.
\end{itemize}

\subsection{MATLAB}
Środowisko MATLAB wykorzystano do przygotowania i~obsługi danych referencyjnych oraz do analizy wyników działania projektu. W szczególności MATLAB służył do:
\begin{itemize}
    \item generowania wektorów testowych wejściowych i~referencyjnych wyników wyjściowych,
    \item przygotowania metadanych testów,
    \item wstępnej analizy porównawczej wyników.
\end{itemize}

MATLAB został wykorzystany również do stworzenia kodu modelującego przetwarzanie strumienia. Kod ten umożliwiał wizualizację kolejnych transformacji wektorów i~weryfikację oczekiwanego formatu danych na wyjściu poszczególnych etapów toru.
